\medskip

\textbf{Les questions 1, 2, 3 et 4 sont indépendantes les unes des autres.}

\medskip

\textbf{Question 1}

\medskip

\textit{Pour cette question, indiquer la lettre de la réponse exacte, en expliquant votre choix.}

On considère la fonction $f$ définie sur $\R$ par 
$f(x) = \e^{- 0,016x} - 2$
et on note $f'$ sa fonction dérivée.

\begin{center}
\begin{tabularx}{\linewidth}{|*{4}{>{\centering \arraybackslash}X|}}\hline
A&B&C&D\\ \hline
\raisebox{-2mm}{$f(0) = - 2$} & 
\raisebox{-2mm}{$f'(x) = \e^{-0,016x}$} & 
$f$ est croissante sur $\R$ & 
$f$ est décroissante sur $\R$\\ \hline
\end{tabularx}
\end{center}

\bigskip

\textbf{Question 2}

\medskip

Soit la fonction $f$ définie sur $[0~;~+ \infty[$ par 
$f(x) = \dfrac{30}{1 + 2x}$.

Est-il vrai que la valeur moyenne de $f$ sur l'intervalle [5~;~15] est supérieure à 3 ?

Justifier la réponse.

\textit{Rappel} : la valeur moyenne de la fonction $f$ sur l'intervalle $[a~;~b]$ est $\displaystyle\dfrac{1}{b - a}\int_a^b f(t)\,\text{d}t$.

\bigskip

\textbf{Question 3}

\medskip

On note $f(t)$ la température (en $\degres C$) d'un café en fonction du temps $t$ (en minute) écoulé depuis sa sortie d'une machine à expresso. À l'instant $t = 0$, la température initiale du café est $83~\degres C$.

On admet que la fonction température est solution sur $[0~;~ +\infty[$ de l'équation différentielle :

\[y' = - 0,08y + 1,84.\]

Déterminer l'expression de l'unique solution $f$ qui vérifie les données précédentes.

\bigskip

\textbf{Question 4}

\medskip

On note $f(t)$ la température (en $\degres C$) d'un café en fonction du temps $t$ (en minute) écoulé depuis sa sortie d'une machine à expresso. On admet, pour tout $t \in [0~;~ +\infty[$, l'expression suivante :

\[f(t) = 60\e^{- 0,08t} + 23.\]

Au bout de combien de temps la température du café sera-t-elle inférieure ou égale à $44~\degres C$ ?

Donner la réponse à la minute près.

\bigskip
